%% ch10 --- Rozciągnij i złóż: dziwne atraktory
%%
%% Zalążek, wygenerowany przez tools/gen_stubs.py z tools/chapters.json.
%% Poniższy opis jest kontraktem tego rozdziału: co ma uzasadnić, co ma dać
%% czytelnikowi i czego NIE ma zabierać sąsiednim rozdziałom. Napisanie
%% rozdziału oznacza usunięcie zalążkowego bloku (tego, który zaczyna się po
%% linii \chapter) i zastąpienie go rozdziałem -- po czym ten plik nie jest
%% już generowany.

\chapter{Rozciągnij i złóż: dziwne atraktory}
\label{ch:strange}

\chapterstub{%
\textit{[Opis do przetłumaczenia --- patrz wydanie angielskie.]}

Argument: every chaotic system does two things over and over \dash{} it
STRETCHES (nearby points separate, the positive Lyapunov exponent) and FOLDS
(the stretched set is put back into a bounded region, or it would fly
apart). A baker kneading dough is the picture, Smale's horseshoe is the
mathematics in words, and the Henon map (1976, a = 1.4, b = 0.3) is a two-
line rule that does it. Each step of the Henon map shrinks area by the
factor |b|, computed as a determinant and checked by measuring, so the
attractor has zero area \dash{} and yet stretching gives it infinite length.
Zooming in repeatedly on the attractor shows the same layered structure at
every scale, which is what makes it strange rather than merely odd, and
hands the geometry to Chapter 11. Payoff: the reader can iterate the Henon
map, show that a cloud of starting points collapses onto the same attractor,
measure the area contraction and the Lyapunov exponent, and explain chaos as
stretch plus fold. Traps to elicit: shrinking area means the motion dies
out; a bounded system cannot be sensitive to its start; zooming in on a
curve always ends with a straight line. Listings: the Henon map; a cloud of
starts collapsing; area contraction measured. Labs: implement one Henon step
and its area factor. Figures: the attractor; two successive zooms; a square
stretched and folded after a few steps.

\medskip
\textbf{Budżet słów:} 3500.
}
